% =====================================================================
%  Anatomía de un LLM — Fase V — Módulo 20
%  Ajuste fino e instrucciones
%  Compilar: latexmk -r ../latexmkrc M20-ajuste-fino.tex
% =====================================================================
\documentclass[11pt,a4paper,twoside,openany]{book}
\usepackage{../estilo-apuntes}

\newcommand{\ApuntesCurso}{Anatomía de un LLM}
\newcommand{\ApuntesModulo}{Módulo 20}

% Notación propia del módulo
\newcommand{\T}{^{\top}}
\newcommand{\Lsft}{\mathcal{L}_{\mathrm{SFT}}}
\newcommand{\rango}{\operatorname{rango}}
\newcommand{\img}{\operatorname{im}}
\newcommand{\Col}{\operatorname{Col}}
\newcommand{\diag}{\operatorname{diag}}
\newcommand{\Fis}{\mathcal{I}}
\newcommand{\tok}[1]{\texttt{\detokenize{#1}}}
\makeatletter\g@addto@macro\UrlBreaks{\do\-}\makeatother

\hypersetup{pdftitle={Módulo 20. Ajuste fino e instrucciones (resumen)},pdfauthor={Anatomía de un LLM}}

\makeatletter
\newlabel{ej:base}{{1.3}{1}{El modelo base no responde}{ejemplo.3}{}}
\newlabel{ej:base@cref}{{[ejemplo][3][]1.3}{[1][1][]1}}
\newlabel{ej:discriminativas}{{1.11}{3}{Tasas discriminativas}{ejemplo.11}{}}
\newlabel{ej:discriminativas@cref}{{[ejemplo][11][]1.11}{[1][3][]3}}
\newlabel{fig:posentrenamiento}{{1.1}{4}{Etapas del entrenamiento de un asistente. El ajuste supervisado (SFT) usa entre miles y cientos de miles de diálogos; puede aplicarse a todos los pesos o, con LoRA, solo a una corrección de bajo rango}{figure.1.1}{}}
\newlabel{fig:posentrenamiento@cref}{{[figure][1][1]1.1}{[1][3][]4}}
\newlabel{eje:1-tasa}{{1.1}{4}{}{ejercicio.1.1}{}}
\newlabel{eje:1-tasa@cref}{{[ejercicio][1][1]1.1}{[1][4][]4}}
\newlabel{eje:1-convexa}{{1.2}{4}{}{ejercicio.1.2}{}}
\newlabel{eje:1-convexa@cref}{{[ejercicio][2][1]1.2}{[1][4][]4}}
\newlabel{eje:1-discriminativas}{{1.3}{4}{}{ejercicio.1.3}{}}
\newlabel{eje:1-discriminativas@cref}{{[ejercicio][3][1]1.3}{[1][4][]4}}
\newlabel{eje:1-stlr}{{1.4}{4}{}{ejercicio.1.4}{}}
\newlabel{eje:1-stlr@cref}{{[ejercicio][4][1]1.4}{[1][4][]4}}
\newlabel{eje:1-media}{{1.5}{4}{}{ejercicio.1.5}{}}
\newlabel{eje:1-media@cref}{{[ejercicio][5][1]1.5}{[1][4][]4}}
\newlabel{ej:mascara}{{2.9}{7}{Con máscara y sin ella}{ejemplo.9}{}}
\newlabel{ej:mascara@cref}{{[ejemplo][9][]2.9}{[1][7][]7}}
\newlabel{fig:plantilla}{{2.1}{8}{Plantilla de chat de un turno (\texttt {us}: usuario; \texttt {as}: asistente) y su máscara. Solo puntúan el contenido de la respuesta y el token de fin que la cierra (sombreados)}{figure.2.1}{}}
\newlabel{fig:plantilla@cref}{{[figure][1][2]2.1}{[1][8][]8}}
\newlabel{eje:2-base}{{2.1}{9}{}{ejercicio.2.1}{}}
\newlabel{eje:2-base@cref}{{[ejercicio][1][2]2.1}{[1][9][]9}}
\newlabel{eje:2-ignore}{{2.2}{9}{}{ejercicio.2.2}{}}
\newlabel{eje:2-ignore@cref}{{[ejercicio][2][2]2.2}{[1][9][]9}}
\newlabel{eje:2-mascara}{{2.3}{9}{}{ejercicio.2.3}{}}
\newlabel{eje:2-mascara@cref}{{[ejercicio][3][2]2.3}{[1][9][]9}}
\newlabel{eje:2-tokenizar}{{2.4}{9}{}{ejercicio.2.4}{}}
\newlabel{eje:2-tokenizar@cref}{{[ejercicio][4][2]2.4}{[1][9][]9}}
\newlabel{eje:2-microlote}{{2.5}{9}{}{ejercicio.2.5}{}}
\newlabel{eje:2-microlote@cref}{{[ejercicio][5][2]2.5}{[1][9][]9}}
\newlabel{eje:2-empaquetar}{{2.6}{10}{}{ejercicio.2.6}{}}
\newlabel{eje:2-empaquetar@cref}{{[ejercicio][6][2]2.6}{[1][10][]10}}
\newlabel{eje:2-multiturno}{{2.7}{10}{}{ejercicio.2.7}{}}
\newlabel{eje:2-multiturno@cref}{{[ejercicio][7][2]2.7}{[1][10][]10}}
\newlabel{ej:olvido-practica}{{3.2}{11}{El olvido en la práctica del módulo}{ejemplo.2}{}}
\newlabel{ej:olvido-practica@cref}{{[ejemplo][2][]3.2}{[1][11][]11}}
\newlabel{ej:muelle}{{3.5}{12}{}{ejemplo.5}{}}
\newlabel{ej:muelle@cref}{{[ejemplo][5][]3.5}{[1][12][]12}}
\newlabel{eje:3-remedios}{{3.1}{14}{}{ejercicio.3.1}{}}
\newlabel{eje:3-remedios@cref}{{[ejercicio][1][3]3.1}{[1][14][]14}}
\newlabel{eje:3-muelle}{{3.2}{14}{}{ejercicio.3.2}{}}
\newlabel{eje:3-muelle@cref}{{[ejercicio][2][3]3.2}{[1][14][]14}}
\newlabel{eje:3-ewc}{{3.3}{14}{}{ejercicio.3.3}{}}
\newlabel{eje:3-ewc@cref}{{[ejercicio][3][3]3.3}{[1][14][]14}}
\newlabel{eje:3-replay}{{3.4}{14}{}{ejercicio.3.4}{}}
\newlabel{eje:3-replay@cref}{{[ejercicio][4][3]3.4}{[1][14][]14}}
\newlabel{eje:3-fisher}{{3.5}{14}{}{ejercicio.3.5}{}}
\newlabel{eje:3-fisher@cref}{{[ejercicio][5][3]3.5}{[1][14][]14}}
\newlabel{eje:3-gbc}{{3.6}{14}{}{ejercicio.3.6}{}}
\newlabel{eje:3-gbc@cref}{{[ejercicio][6][3]3.6}{[1][14][]14}}
\newlabel{fig:adaptador}{{4.1}{17}{Adaptador con cuello de botella. Si $W_{\uparrow }$ y $b_{\uparrow }$ son nulos, el adaptador es la identidad}{figure.4.1}{}}
\newlabel{fig:adaptador@cref}{{[figure][1][4]4.1}{[1][17][]17}}
\newlabel{eje:4-congelar}{{4.1}{19}{}{ejercicio.4.1}{}}
\newlabel{eje:4-congelar@cref}{{[ejercicio][1][4]4.1}{[1][19][]19}}
\newlabel{eje:4-adaptadores}{{4.2}{19}{}{ejercicio.4.2}{}}
\newlabel{eje:4-adaptadores@cref}{{[ejercicio][2][4]4.2}{[1][19][]19}}
\newlabel{eje:4-prefijo}{{4.3}{20}{}{ejercicio.4.3}{}}
\newlabel{eje:4-prefijo@cref}{{[ejercicio][3][4]4.3}{[1][19][]20}}
\newlabel{eje:4-subespacio}{{4.4}{20}{}{ejercicio.4.4}{}}
\newlabel{eje:4-subespacio@cref}{{[ejercicio][4][4]4.4}{[1][20][]20}}
\newlabel{eje:4-proyeccion}{{4.5}{20}{}{ejercicio.4.5}{}}
\newlabel{eje:4-proyeccion@cref}{{[ejercicio][5][4]4.5}{[1][20][]20}}
\newlabel{fig:lora}{{5.1}{21}{Capa lineal con LoRA. La entrada atraviesa la matriz congelada y, en paralelo, el cuello de dimensión $r$ formado por $A$ y $B$. Como $B$ empieza en cero, al principio $h=Wx$}{figure.5.1}{}}
\newlabel{fig:lora@cref}{{[figure][1][5]5.1}{[1][21][]21}}
\newlabel{ej:lora-mano}{{5.3}{21}{Una capa LoRA a mano}{ejemplo.3}{}}
\newlabel{ej:lora-mano@cref}{{[ejemplo][3][]5.3}{[1][21][]21}}
\newlabel{ej:energia}{{5.13}{24}{}{ejemplo.13}{}}
\newlabel{ej:energia@cref}{{[ejemplo][13][]5.13}{[1][24][]24}}
\newlabel{eje:5-vf}{{5.1}{24}{}{ejercicio.5.1}{}}
\newlabel{eje:5-vf@cref}{{[ejercicio][1][5]5.1}{[1][24][]24}}
\newlabel{eje:5-contar}{{5.2}{25}{}{ejercicio.5.2}{}}
\newlabel{eje:5-contar@cref}{{[ejercicio][2][5]5.2}{[1][24][]25}}
\newlabel{eje:5-fusion}{{5.3}{25}{}{ejercicio.5.3}{}}
\newlabel{eje:5-fusion@cref}{{[ejercicio][3][5]5.3}{[1][25][]25}}
\newlabel{eje:5-grad}{{5.4}{25}{}{ejercicio.5.4}{}}
\newlabel{eje:5-grad@cref}{{[ejercicio][4][5]5.4}{[1][25][]25}}
\newlabel{eje:5-svd}{{5.5}{25}{}{ejercicio.5.5}{}}
\newlabel{eje:5-svd@cref}{{[ejercicio][5][5]5.5}{[1][25][]25}}
\newlabel{eje:5-rango}{{5.6}{25}{}{ejercicio.5.6}{}}
\newlabel{eje:5-rango@cref}{{[ejercicio][6][5]5.6}{[1][25][]25}}
\newlabel{eje:5-init}{{5.7}{25}{}{ejercicio.5.7}{}}
\newlabel{eje:5-init@cref}{{[ejercicio][7][5]5.7}{[1][25][]25}}
\newlabel{ej:cuantiles2}{{6.4}{27}{Una versión simplificada con 2 bits}{ejemplo.4}{}}
\newlabel{ej:cuantiles2@cref}{{[ejemplo][4][]6.4}{[1][27][]27}}
\newlabel{fig:decision}{{6.1}{29}{Diagrama de decisión orientativo para adaptar un modelo de lenguaje. Las preguntas (recuadros rectos) se recorren de arriba abajo; las respuestas están sombreadas}{figure.6.1}{}}
\newlabel{fig:decision@cref}{{[figure][1][6]6.1}{[1][28][]29}}
\newlabel{eje:6-rag}{{6.1}{29}{}{ejercicio.6.1}{}}
\newlabel{eje:6-rag@cref}{{[ejercicio][1][6]6.1}{[1][29][]29}}
\newlabel{eje:6-cuantiles}{{6.2}{29}{}{ejercicio.6.2}{}}
\newlabel{eje:6-cuantiles@cref}{{[ejercicio][2][6]6.2}{[1][29][]29}}
\newlabel{eje:6-doble}{{6.3}{29}{}{ejercicio.6.3}{}}
\newlabel{eje:6-doble@cref}{{[ejercicio][3][6]6.3}{[1][29][]29}}
\newlabel{eje:6-memoria}{{6.4}{30}{}{ejercicio.6.4}{}}
\newlabel{eje:6-memoria@cref}{{[ejercicio][4][6]6.4}{[1][30][]30}}
\newlabel{eje:6-contabilidad}{{6.5}{30}{}{ejercicio.6.5}{}}
\newlabel{eje:6-contabilidad@cref}{{[ejercicio][5][6]6.5}{[1][30][]30}}
\newlabel{eje:6-rango-limite}{{6.6}{30}{}{ejercicio.6.6}{}}
\newlabel{eje:6-rango-limite@cref}{{[ejercicio][6][6]6.6}{[1][30][]30}}
\newlabel{ap:soluciones}{{A}{31}{Soluciones de los ejercicios}{appendix.A}{}}
\newlabel{ap:soluciones@cref}{{[appendix][1][2147483647]A}{[1][31][]31}}
\newlabel{sol:eje:1-tasa}{{A.1}{31}{Tema 1}{section*.3}{}}
\newlabel{sol:eje:1-tasa@cref}{{[subappendix][1][2147483647,1]A.1}{[1][31][]31}}
\newlabel{sol:eje:1-convexa}{{A.1}{31}{Tema 1}{section*.4}{}}
\newlabel{sol:eje:1-convexa@cref}{{[subappendix][1][2147483647,1]A.1}{[1][31][]31}}
\newlabel{sol:eje:1-discriminativas}{{A.1}{32}{Tema 1}{section*.5}{}}
\newlabel{sol:eje:1-discriminativas@cref}{{[subappendix][1][2147483647,1]A.1}{[1][32][]32}}
\newlabel{sol:eje:1-stlr}{{A.1}{32}{Tema 1}{section*.6}{}}
\newlabel{sol:eje:1-stlr@cref}{{[subappendix][1][2147483647,1]A.1}{[1][32][]32}}
\newlabel{sol:eje:1-media}{{A.1}{32}{Tema 1}{section*.7}{}}
\newlabel{sol:eje:1-media@cref}{{[subappendix][1][2147483647,1]A.1}{[1][32][]32}}
\newlabel{sol:eje:2-base}{{A.2}{33}{Tema 2}{section*.8}{}}
\newlabel{sol:eje:2-base@cref}{{[subappendix][2][2147483647,1]A.2}{[1][33][]33}}
\newlabel{sol:eje:2-ignore}{{A.2}{33}{Tema 2}{section*.9}{}}
\newlabel{sol:eje:2-ignore@cref}{{[subappendix][2][2147483647,1]A.2}{[1][33][]33}}
\newlabel{sol:eje:2-mascara}{{A.2}{34}{Tema 2}{section*.10}{}}
\newlabel{sol:eje:2-mascara@cref}{{[subappendix][2][2147483647,1]A.2}{[1][34][]34}}
\newlabel{sol:eje:2-tokenizar}{{A.2}{34}{Tema 2}{section*.11}{}}
\newlabel{sol:eje:2-tokenizar@cref}{{[subappendix][2][2147483647,1]A.2}{[1][34][]34}}
\newlabel{sol:eje:2-microlote}{{A.2}{35}{Tema 2}{section*.12}{}}
\newlabel{sol:eje:2-microlote@cref}{{[subappendix][2][2147483647,1]A.2}{[1][34][]35}}
\newlabel{sol:eje:2-empaquetar}{{A.2}{35}{Tema 2}{section*.13}{}}
\newlabel{sol:eje:2-empaquetar@cref}{{[subappendix][2][2147483647,1]A.2}{[1][35][]35}}
\newlabel{sol:eje:2-multiturno}{{A.2}{35}{Tema 2}{section*.14}{}}
\newlabel{sol:eje:2-multiturno@cref}{{[subappendix][2][2147483647,1]A.2}{[1][35][]35}}
\newlabel{sol:eje:3-remedios}{{A.3}{36}{Tema 3}{section*.15}{}}
\newlabel{sol:eje:3-remedios@cref}{{[subappendix][3][2147483647,1]A.3}{[1][36][]36}}
\newlabel{sol:eje:3-muelle}{{A.3}{36}{Tema 3}{section*.16}{}}
\newlabel{sol:eje:3-muelle@cref}{{[subappendix][3][2147483647,1]A.3}{[1][36][]36}}
\newlabel{sol:eje:3-ewc}{{A.3}{37}{Tema 3}{section*.17}{}}
\newlabel{sol:eje:3-ewc@cref}{{[subappendix][3][2147483647,1]A.3}{[1][37][]37}}
\newlabel{sol:eje:3-replay}{{A.3}{37}{Tema 3}{section*.18}{}}
\newlabel{sol:eje:3-replay@cref}{{[subappendix][3][2147483647,1]A.3}{[1][37][]37}}
\newlabel{sol:eje:3-fisher}{{A.3}{37}{Tema 3}{section*.19}{}}
\newlabel{sol:eje:3-fisher@cref}{{[subappendix][3][2147483647,1]A.3}{[1][37][]37}}
\newlabel{sol:eje:3-gbc}{{A.3}{38}{Tema 3}{section*.20}{}}
\newlabel{sol:eje:3-gbc@cref}{{[subappendix][3][2147483647,1]A.3}{[1][38][]38}}
\newlabel{sol:eje:4-congelar}{{A.4}{38}{Tema 4}{section*.21}{}}
\newlabel{sol:eje:4-congelar@cref}{{[subappendix][4][2147483647,1]A.4}{[1][38][]38}}
\newlabel{sol:eje:4-adaptadores}{{A.4}{39}{Tema 4}{section*.22}{}}
\newlabel{sol:eje:4-adaptadores@cref}{{[subappendix][4][2147483647,1]A.4}{[1][39][]39}}
\newlabel{sol:eje:4-prefijo}{{A.4}{39}{Tema 4}{section*.23}{}}
\newlabel{sol:eje:4-prefijo@cref}{{[subappendix][4][2147483647,1]A.4}{[1][39][]39}}
\newlabel{sol:eje:4-subespacio}{{A.4}{39}{Tema 4}{section*.24}{}}
\newlabel{sol:eje:4-subespacio@cref}{{[subappendix][4][2147483647,1]A.4}{[1][39][]39}}
\newlabel{sol:eje:4-proyeccion}{{A.4}{40}{Tema 4}{section*.25}{}}
\newlabel{sol:eje:4-proyeccion@cref}{{[subappendix][4][2147483647,1]A.4}{[1][40][]40}}
\newlabel{sol:eje:5-vf}{{A.5}{40}{Tema 5}{section*.26}{}}
\newlabel{sol:eje:5-vf@cref}{{[subappendix][5][2147483647,1]A.5}{[1][40][]40}}
\newlabel{sol:eje:5-contar}{{A.5}{41}{Tema 5}{section*.27}{}}
\newlabel{sol:eje:5-contar@cref}{{[subappendix][5][2147483647,1]A.5}{[1][41][]41}}
\newlabel{sol:eje:5-fusion}{{A.5}{41}{Tema 5}{section*.28}{}}
\newlabel{sol:eje:5-fusion@cref}{{[subappendix][5][2147483647,1]A.5}{[1][41][]41}}
\newlabel{sol:eje:5-grad}{{A.5}{42}{Tema 5}{section*.29}{}}
\newlabel{sol:eje:5-grad@cref}{{[subappendix][5][2147483647,1]A.5}{[1][41][]42}}
\newlabel{sol:eje:5-svd}{{A.5}{42}{Tema 5}{section*.30}{}}
\newlabel{sol:eje:5-svd@cref}{{[subappendix][5][2147483647,1]A.5}{[1][42][]42}}
\newlabel{sol:eje:5-rango}{{A.5}{42}{Tema 5}{section*.31}{}}
\newlabel{sol:eje:5-rango@cref}{{[subappendix][5][2147483647,1]A.5}{[1][42][]42}}
\newlabel{sol:eje:5-init}{{A.5}{43}{Tema 5}{section*.32}{}}
\newlabel{sol:eje:5-init@cref}{{[subappendix][5][2147483647,1]A.5}{[1][43][]43}}
\newlabel{sol:eje:6-rag}{{A.6}{43}{Tema 6}{section*.33}{}}
\newlabel{sol:eje:6-rag@cref}{{[subappendix][6][2147483647,1]A.6}{[1][43][]43}}
\newlabel{sol:eje:6-cuantiles}{{A.6}{44}{Tema 6}{section*.34}{}}
\newlabel{sol:eje:6-cuantiles@cref}{{[subappendix][6][2147483647,1]A.6}{[1][44][]44}}
\newlabel{sol:eje:6-doble}{{A.6}{44}{Tema 6}{section*.35}{}}
\newlabel{sol:eje:6-doble@cref}{{[subappendix][6][2147483647,1]A.6}{[1][44][]44}}
\newlabel{sol:eje:6-memoria}{{A.6}{45}{Tema 6}{section*.36}{}}
\newlabel{sol:eje:6-memoria@cref}{{[subappendix][6][2147483647,1]A.6}{[1][44][]45}}
\newlabel{sol:eje:6-contabilidad}{{A.6}{45}{Tema 6}{section*.37}{}}
\newlabel{sol:eje:6-contabilidad@cref}{{[subappendix][6][2147483647,1]A.6}{[1][45][]45}}
\newlabel{sol:eje:6-rango-limite}{{A.6}{45}{Tema 6}{section*.38}{}}
\newlabel{sol:eje:6-rango-limite@cref}{{[subappendix][6][2147483647,1]A.6}{[1][45][]45}}
\makeatother

\begin{document}
\frontmatter
\begin{titlepage}
\centering
\vspace*{2.5cm}
{\large\scshape Anatomía de un LLM}\par
\vspace{0.4cm}
{\normalsize Fundamentos matemáticos de los modelos de lenguaje}\par
\vspace{2.2cm}
\rule{\textwidth}{0.8pt}\par\vspace{0.5cm}
{\small\scshape Fase V \textperiodcentered\ Alineamiento, aprendizaje por refuerzo, agentes e interpretabilidad}\par\vspace{0.3cm}
{\Huge\bfseries Módulo 20\par}\vspace{0.3cm}
{\LARGE Ajuste fino e instrucciones\par}
\vspace{0.2cm}
{\large Ajuste supervisado, olvido catastrófico y adaptación eficiente en parámetros\par}
\vspace{0.5cm}\rule{\textwidth}{0.4pt}\par
\vspace{1.6cm}
{\large Resumen del módulo}\par
\vfill
{\small Edición de octubre de 2026}\par
\end{titlepage}

\chapter*{Sobre este resumen}
\addcontentsline{toc}{chapter}{Sobre este resumen}
\markboth{Sobre este resumen}{}

Este documento reúne lo indispensable de los apuntes del Módulo 20: el texto expositivo, todas las definiciones y todos los resultados (teoremas, proposiciones, lemas y corolarios), las tablas y las observaciones de contenido. Se han suprimido las demostraciones, los ejemplos, las lecturas, las figuras, los ejercicios y sus soluciones. La numeración de temas, secciones, enunciados y ecuaciones es la misma que en los apuntes, de modo que cualquier referencia puede buscarse allí; las que remiten a elementos suprimidos (un ejemplo, una figura, una ecuación de una demostración) conservan su número en los apuntes.

\section*{Notación}
\begin{center}\small
\begin{tabularx}{\textwidth}{@{}l>{\raggedright\arraybackslash}X@{}}
\toprule
Símbolo & Significado \\
\midrule
$\theta_0$, $\theta$ & parámetros del modelo preentrenado y parámetros ajustados \\
$u=u_{1:n}$, $y=y_{1:m}$ & instrucción (\emph{prompt}) y respuesta, como secuencias de tokens \\
$x_{1:T}$, $m_t\in\{0,1\}$ & secuencia completa (instrucción seguida de respuesta) y máscara de las posiciones que puntúan \\
$\Lsft$ & pérdida del ajuste supervisado (entropía cruzada enmascarada) \\
$\eta$, $\eta^{l}$ & tasa de aprendizaje global y de la capa $l$ \\
$\lambda$ & intensidad de la penalización que ancla los parámetros a $\theta_0$ \\
$\Fis(\theta)$, $F_i$ & matriz de información de Fisher y su diagonal \\
$D$, $d'$ & número total de parámetros y dimensión del subespacio entrenable (tema~4) \\
$W\in\R^{d\times k}$, $x\in\R^{k}$ & en los temas~4 a~6, matriz preentrenada congelada y entrada de la capa \\
$A\in\R^{r\times k}$, $B\in\R^{d\times r}$ & factores de LoRA; $r$ es el rango \\
$\alpha$, $s=\alpha/r$ & hiperparámetro de escala de LoRA y factor efectivo \\
$\sigma_1\ge\sigma_2\ge\cdots$ & valores singulares \\
$\Phi$, $\Phi^{-1}$ & función de distribución de la normal estándar y su inversa (función cuantil) \\
$N$, $N_{\mathrm{LoRA}}$ & parámetros del modelo base y de los adaptadores \\
\bottomrule
\end{tabularx}
\end{center}


\tableofcontents

\mainmatter
% =====================================================================
\chapter{Transferencia y ajuste fino}
\label{tema:1}
% =====================================================================

Este tema formula el ajuste fino como un problema de optimización que parte de un modelo ya entrenado y presenta las técnicas clásicas de transferencia que preceden al posentrenamiento de los modelos de lenguaje actuales. La distinción entre extracción de características y ajuste completo sigue el §14.2 de Zhang \emph{et al.}~\cite{d2l}; las técnicas de ULMFiT se toman del artículo de Howard y Ruder~\cite[§3.2 y~§3.3]{howard-ruder}, y la visión del posentrenamiento, del capítulo~8 de Jurafsky y Martin~\cite{jm3} y de la clase~8 de CS224N~\cite{cs224n}.

\section{Modelo base, tarea de origen y tarea de destino}

\begin{definicion}[Aprendizaje por transferencia]\label{def:transferencia}
Sea $f_\theta$ una familia de modelos con parámetros $\theta\in\R^{D}$. Se llama \emph{tarea de origen} a la minimización de una pérdida $\mathcal{L}_{\mathrm{orig}}$ sobre un conjunto de datos grande, y \emph{modelo preentrenado} a unos parámetros $\theta_0$ que la minimizan aproximadamente. Dada una \emph{tarea de destino} con datos $\mathcal{D}$ y pérdida $\mathcal{L}_{\mathcal{D}}$, el \emph{ajuste fino} consiste en minimizar aproximadamente $\mathcal{L}_{\mathcal{D}}$ mediante un método iterativo de primer orden inicializado en $\theta_0$, normalmente con pocos pasos y una tasa de aprendizaje pequeña.
\end{definicion}

En los modelos de lenguaje, la tarea de origen es la predicción del token siguiente sobre un corpus masivo (módulos~15 y~16), y $\theta_0$ define un \emph{modelo base}.

\begin{definicion}[Modelo base]\label{def:base}
Un \emph{modelo base} es un modelo de lenguaje autorregresivo entrenado únicamente con la entropía cruzada del token siguiente sobre texto general. Para un contexto $x_{1:t-1}$, devuelve una aproximación de la distribución condicional del corpus, $p_{\theta_0}(x_t\mid x_{1:t-1})$.
\end{definicion}

\addtocounter{definicion}{1}%

\begin{observacion}\label{obs:implicita}
Como el problema de destino no es convexo y la optimización es corta, el resultado depende de $\theta_0$ no solo como punto de partida: el método iterativo permanece cerca de $\theta_0$ y actúa como una regularización implícita. El \cref{tema:3} hace explícita esta regularización mediante una penalización de la distancia a $\theta_0$.
\end{observacion}

\section{Extracción de características y ajuste completo}

Sea $\phi_\theta(x)\in\R^{h}$ la representación que la red asigna a una entrada (por ejemplo, el vector oculto de la última capa) y sea una \emph{cabeza} lineal $z=V\phi_\theta(x)+b$, con $V\in\R^{K\times h}$, que produce los logits de $K$ clases.

\begin{definicion}[Extracción de características y ajuste completo]\label{def:sonda}
En la \emph{extracción de características} (o \emph{sonda lineal}) se fijan los parámetros del cuerpo en $\theta=\theta_0$ y solo se entrenan $(V,b)$. En el \emph{ajuste completo} se entrenan a la vez $\theta$, inicializado en $\theta_0$, y la cabeza, inicializada al azar.
\end{definicion}

\begin{proposicion}[La sonda lineal es convexa]\label{prop:sonda-convexa}
Con el cuerpo congelado, la pérdida de entropía cruzada media
\[
J(V,b)=\frac{1}{n}\sum_{i=1}^{n}\Bigl(\LSE\bigl(V\phi_{\theta_0}(x_i)+b\bigr)-\bigl(V\phi_{\theta_0}(x_i)+b\bigr)_{y_i}\Bigr)
\]
es una función convexa de $(V,b)$.
\end{proposicion}

En el ajuste completo, $\phi_\theta$ depende de $\theta$ de forma no lineal y la convexidad se pierde. A cambio, el cuerpo puede adaptar sus representaciones a la tarea. La receta habitual~\cite[§14.2]{d2l} usa una tasa de aprendizaje pequeña en el cuerpo, que ya está cerca de una buena solución, y una mayor (por ejemplo, diez veces mayor) en la cabeza, que parte de cero.

\section{ULMFiT: ajuste fino de un modelo de lenguaje}

Howard y Ruder~\cite{howard-ruder} propusieron en 2018 el primer procedimiento general para ajustar un modelo de lenguaje preentrenado a tareas de clasificación de textos. Consta de tres etapas: preentrenamiento del modelo de lenguaje en un corpus general, ajuste del modelo de lenguaje al texto del dominio de destino y ajuste del clasificador. Introduce tres técnicas para que el ajuste no destruya lo aprendido.

\begin{definicion}[Tasas de aprendizaje discriminativas]\label{def:discriminativas}
Sea una red de $L$ capas con parámetros $\theta^{1},\dots,\theta^{L}$. El \emph{ajuste discriminativo} asigna a cada capa su propia tasa, $\theta^{l}_t=\theta^{l}_{t-1}-\eta^{l}\nabla_{\theta^{l}}J$, con $\eta^{l-1}=\eta^{l}/2{,}6$: cada capa aprende $2{,}6$ veces más despacio que la siguiente. La tasa $\eta^{L}$ de la última capa se elige ajustando solo esa capa~\cite[§3.2]{howard-ruder}.
\end{definicion}

Por inducción, $\eta^{l}=\eta^{L}/2{,}6^{\,L-l}$. La justificación es que las capas inferiores capturan información más general, que conviene modificar menos.

\begin{definicion}[Tasa triangular inclinada]\label{def:stlr}
Sean $T$ el número total de iteraciones, $\mathrm{cf}\in(0,1)$ la fracción de calentamiento, $\rho>1$ el cociente entre la tasa máxima y la mínima y $\eta_{\max}$ la tasa máxima. Con $c=\lfloor T\,\mathrm{cf}\rfloor$, la \emph{tasa triangular inclinada} (STLR) en la iteración $t$ es
\[
p_t=\begin{cases} t/c, & t<c,\\[2pt] 1-\dfrac{t-c}{c\,(1/\mathrm{cf}-1)}, & t\ge c,\end{cases}\qquad
\eta_t=\eta_{\max}\,\frac{1+p_t(\rho-1)}{\rho}.
\]
Los valores por defecto del artículo son $\mathrm{cf}=0{,}1$, $\rho=32$ y $\eta_{\max}=0{,}01$~\cite[§3.2]{howard-ruder}.
\end{definicion}

\begin{proposicion}[Forma de la tasa triangular]\label{prop:stlr}
Supóngase que $T\,\mathrm{cf}$ es entero, de modo que $c=T\,\mathrm{cf}$. Entonces:
\begin{enumerate}
  \item $\eta_t$ es una función afín de $p_t$, creciente, con $\eta=\eta_{\max}/\rho$ si $p_t=0$ y $\eta=\eta_{\max}$ si $p_t=1$;
  \item $p_t$ crece linealmente de $0$ a $1$ en $[0,c]$ y decrece linealmente de $1$ a $0$ en $[c,T]$;
  \item en consecuencia, la tasa sube en $c$ iteraciones desde $\eta_{\max}/\rho$ hasta $\eta_{\max}$ y baja en $T-c$ iteraciones hasta $\eta_{\max}/\rho$.
\end{enumerate}
\end{proposicion}

\begin{definicion}[Descongelación gradual]\label{def:descongelacion}
En la \emph{descongelación gradual} se entrena primero solo la última capa durante una época; después se descongela la capa anterior y se entrenan todas las descongeladas otra época, y así sucesivamente hasta entrenar todas las capas~\cite[§3.3]{howard-ruder}.
\end{definicion}

\addtocounter{definicion}{1}%

\begin{observacion}
Las tres técnicas persiguen el mismo fin: evitar que los primeros pasos del ajuste, con gradientes grandes procedentes de una cabeza aleatoria, destruyan las representaciones del modelo preentrenado. Howard y Ruder las presentan como remedio al \emph{olvido catastrófico}, que el \cref{tema:3} estudia con detalle. La tasa triangular es, además, un antecedente directo del calentamiento seguido de decaimiento que se usa en el entrenamiento de los Transformer (módulo~7).
\end{observacion}

\section{Del ajuste por tareas al posentrenamiento}

ULMFiT añade una cabeza de clasificación y produce un modelo para una sola tarea. En los modelos de lenguaje generativos se conserva la cabeza de lenguaje: la tarea de destino se formula también como predicción del token siguiente, pero sobre textos que tienen la forma de una conversación entre un usuario y un asistente. Jurafsky y Martin~\cite[cap.~8]{jm3} llaman \emph{posentrenamiento} al conjunto de etapas que siguen al preentrenamiento y lo dividen en el ajuste con instrucciones (§8.1, \cref{tema:2}) y el aprendizaje a partir de preferencias (§8.3 y~§8.4, módulo~21). La receta en tres pasos de InstructGPT~\cite{ouyang} (ajuste supervisado, modelo de recompensa y aprendizaje por refuerzo) es la referencia publicada que siguen casi todos los asistentes (\cref{fig:posentrenamiento}).

\addtocounter{figure}{1}%

\begin{observacion}[Información pública sobre modelos comerciales]
Anthropic, Google y xAI no publican los datos ni los hiperparámetros del ajuste supervisado de Claude, Gemini y Grok. Lo público es el esquema general. Anthropic describió en 2022 el método \emph{Constitutional AI}~\cite{bai-cai}, que incluye una fase supervisada en la que el modelo critica y revisa sus propias respuestas según una lista de principios y se ajusta sobre las respuestas revisadas, seguida de aprendizaje por refuerzo con retroalimentación de un modelo (módulo~21). El informe técnico de Gemini~\cite[§6.3]{gemini} describe, a grandes rasgos, un ajuste supervisado seguido de aprendizaje por refuerzo con retroalimentación humana. xAI no ha publicado una descripción equivalente del posentrenamiento de Grok.
\end{observacion}


% =====================================================================
\chapter{Ajuste supervisado con instrucciones}
\label{tema:2}
% =====================================================================

Este tema define el ajuste supervisado (SFT, del inglés \emph{supervised fine-tuning}) de un modelo de lenguaje: la misma entropía cruzada del preentrenamiento, evaluada solo en las posiciones de la respuesta. Se estudian los datos de instrucciones, la pérdida enmascarada y su interpretación probabilística, los detalles de construcción de lotes y la plantilla de chat. Se sigue el §8.1 de Jurafsky y Martin~\cite{jm3}, la clase~15 de CS336~\cite{cs336} y su tarea~5~\cite[§4.2]{cs336a5}, y los artículos de FLAN~\cite{wei-flan}, T0~\cite{sanh-t0}, InstructGPT~\cite[§3.2 y~§3.5]{ouyang} y LIMA~\cite{zhou-lima}.

\section{Datos de instrucciones}

\begin{definicion}[Conjunto de instrucciones]\label{def:instrucciones}
Un \emph{conjunto de instrucciones} es una colección de pares $\{(u^{(i)},y^{(i)})\}_{i=1}^{n}$ en la que $u^{(i)}$ es una instrucción en lenguaje natural, que puede incluir un texto de entrada, e $y^{(i)}$ es la respuesta deseada. El \emph{ajuste con instrucciones} (\emph{instruction tuning}) es el ajuste supervisado sobre un conjunto de instrucciones que abarca muchas tareas distintas.
\end{definicion}

Los conjuntos se construyen de dos maneras~\cite[§8.1.1]{jm3}: reformulando como instrucciones conjuntos de datos supervisados ya existentes, mediante plantillas, o recogiendo demostraciones escritas por personas (o generadas por otros modelos). La colección Flan 2022, por ejemplo, reúne 15 millones de ejemplos de 1836 tareas~\cite[§8.1.1]{jm3}.

\begin{definicion}[Generalización a tareas no vistas]\label{def:no-vistas}
Se agrupan las tareas en \emph{grupos} (traducción, resumen, inferencia, preguntas y respuestas, etc.). Para medir la generalización a tareas no vistas se ajusta el modelo con todos los grupos salvo uno y se evalúa sin ejemplos (\emph{zero-shot}) en las tareas del grupo excluido.
\end{definicion}

Con este protocolo, Wei \emph{et al.}~\cite{wei-flan} ajustaron un modelo de 137\,000 millones de parámetros con más de 60 conjuntos de datos expresados mediante plantillas de instrucciones (FLAN). El modelo ajustado superó, sin ejemplos, a GPT-3 de 175\,000 millones de parámetros en 20 de los 25 conjuntos evaluados, y el beneficio dependía de la escala: en los modelos más pequeños estudiados, el ajuste con instrucciones empeoraba el resultado en las tareas no vistas. Sanh \emph{et al.}~\cite{sanh-t0} obtuvieron resultados análogos con T0, un modelo de tipo T5 de 11\,000 millones de parámetros ajustado con una colección de conjuntos de datos reformulados mediante múltiples \emph{prompts}, que igualaba o superaba en varias pruebas a modelos hasta 16 veces mayores.

\begin{observacion}[El ajuste supervisado en InstructGPT]\label{obs:instructgpt}
En InstructGPT~\cite{ouyang}, el primer paso usa unas 13\,000 instrucciones de entrenamiento con demostraciones escritas por anotadores~\cite[§3.2]{ouyang}. El modelo se ajusta durante 16 épocas con decaimiento coseno de la tasa de aprendizaje y \emph{dropout} residual de $0{,}2$; la pérdida de validación empieza a subir después de la primera época, pero entrenar más épocas mejora la puntuación del modelo de recompensa y las preferencias humanas~\cite[§3.5]{ouyang}. La pérdida de validación del SFT no es, por tanto, un indicador fiable de la calidad final.
\end{observacion}

\section{La pérdida enmascarada}

Sean $u=u_{1:n}$ la instrucción e $y=y_{1:m}$ la respuesta, ambas ya tokenizadas, y sea $x_{1:T}=(u_1,\dots,u_n,y_1,\dots,y_m)$, con $T=n+m$, la secuencia concatenada. El último token de $y$ es el token de fin de turno.

\begin{definicion}[Pérdida del ajuste supervisado]\label{def:sft}
Sea $m_t\in\{0,1\}$ una \emph{máscara} que vale $1$ en las posiciones de la respuesta ($t>n$) y $0$ en las de la instrucción. La \emph{pérdida del ajuste supervisado} es la entropía cruzada enmascarada
\[
\Lsft(\theta)=-\frac{1}{\sum_{t=1}^{T}m_t}\sum_{t=1}^{T}m_t\ln p_\theta(x_t\mid x_{1:t-1}).
\]
\end{definicion}

Se divide por el número de posiciones que cuentan, no por $T$. Si $m_t=1$ para todo $t$, se recupera la pérdida del preentrenamiento.

\begin{proposicion}[Interpretación probabilística]\label{prop:sft-condicional}
Con la máscara de la \cref{def:sft},
\[
\Lsft(\theta)=-\frac{1}{m}\ln p_\theta(y\mid u),\qquad\text{donde}\quad p_\theta(y\mid u)=\prod_{j=1}^{m}p_\theta(y_j\mid u,y_{1:j-1}).
\]
En consecuencia, minimizar la suma de $m^{(i)}\Lsft^{(i)}$ sobre un conjunto de instrucciones equivale a estimar por máxima verosimilitud la distribución condicional de la respuesta dada la instrucción.
\end{proposicion}

\begin{observacion}
Sin máscara se minimizaría $-\frac{1}{T}\ln p_\theta(u,y)=-\frac1T\bigl(\ln p_\theta(u)+\ln p_\theta(y\mid u)\bigr)$: el modelo dedicaría parte de su capacidad a imitar las instrucciones de los usuarios, que no tiene que generar. Este es el sentido de la indicación de Jurafsky y Martin de entrenar solo sobre la respuesta~\cite[§8.1]{jm3}.
\end{observacion}

\begin{proposicion}[Gradiente en las posiciones enmascaradas]\label{prop:grad-mascara}
Sea $z_t\in\R^{K}$ el vector de logits de la posición $t$ y $p_t=\softmax(z_t)$. Entonces
\[
\frac{\partial\Lsft}{\partial z_t}=\frac{m_t}{\sum_{s}m_s}\,\bigl(p_t-e_{x_t}\bigr),
\]
donde $e_{x_t}$ es el vector canónico del token objetivo. En particular, el gradiente respecto a los logits de una posición enmascarada es nulo.
\end{proposicion}

\begin{observacion}
Que el gradiente respecto a los logits de la instrucción sea nulo no implica que la instrucción no influya en el aprendizaje: las posiciones de la respuesta atienden a las de la instrucción, de modo que el gradiente llega a los estados ocultos de la instrucción a través de la atención. Lo que se suprime es la señal que obligaría a \emph{predecir} los tokens de la instrucción.
\end{observacion}

\addtocounter{definicion}{1}%

\begin{observacion}[Implementación]
En PyTorch no se multiplica por la máscara de forma explícita: se asigna la etiqueta $-100$ a las posiciones que no cuentan y la función \texttt{cross\_\allowbreak entropy} con \texttt{ignore\_index=-100}, el valor por defecto, las excluye del numerador y del denominador de la media.
\end{observacion}

\section{Lotes: relleno, normalización, empaquetado y acumulación}

Las secuencias de un lote tienen longitudes distintas. Hay dos formas de formar tensores rectangulares.

\begin{definicion}[Relleno y empaquetado]\label{def:relleno}
En el \emph{relleno} (\emph{padding}) cada secuencia se completa con un token de relleno hasta la longitud máxima del lote, y la máscara vale $0$ en las posiciones de relleno. En el \emph{empaquetado} (\emph{packing}) se concatenan los ejemplos, separados por el token de fin, y la cadena resultante se corta en trozos consecutivos de longitud fija $L$.
\end{definicion}

El empaquetado no desperdicia posiciones, pero un trozo puede contener el final de un ejemplo y el principio de otro; la máscara causal permite entonces que un ejemplo atienda al anterior, lo que suele tolerarse porque el token de fin los separa.

Sea un lote de $B$ secuencias con máscaras $m^{(i)}_t$ y pérdidas por token $\ell^{(i)}_t=-\ln p_\theta(x^{(i)}_t\mid x^{(i)}_{1:t-1})$, y sea $M_i=\sum_t m^{(i)}_t$. Hay dos normalizaciones habituales:
\[
J_{\mathrm{tok}}=\frac{\sum_{i}\sum_t m^{(i)}_t\ell^{(i)}_t}{\sum_i M_i},\qquad
J_{\mathrm{sec}}=\frac{1}{B}\sum_{i=1}^{B}\frac{1}{c}\sum_t m^{(i)}_t\ell^{(i)}_t,
\]
con $c>0$ una constante fija (por ejemplo, la longitud máxima). En $J_{\mathrm{tok}}$ cada secuencia pesa en proporción a la longitud de su respuesta; en $J_{\mathrm{sec}}$ todas las secuencias pesan lo mismo por cada token, independientemente del lote en el que caigan. La tarea~5 de CS336 implementa una normalización del segundo tipo: una suma enmascarada dividida por una constante~\cite[§4.2]{cs336a5}.

\begin{proposicion}[Acumulación de gradientes]\label{prop:acumulacion}
Divídase un lote de $B$ secuencias en $G$ microlotes $\mathcal{B}_1,\dots,\mathcal{B}_G$ de $B/G$ secuencias. Sea $J_g$ la pérdida $J_{\mathrm{sec}}$ del microlote $g$ calculada con tamaño $B/G$. Si cada microlote retropropaga $J_g/G$ y se suman los gradientes antes de actualizar, el gradiente acumulado coincide con el de $J_{\mathrm{sec}}$ sobre el lote completo.
\end{proposicion}

\begin{observacion}
Con la normalización $J_{\mathrm{tok}}$ la equivalencia falla en general: cada microlote divide por su propio número de tokens de respuesta y la suma de los cocientes no es el cociente de las sumas, salvo que todos los microlotes tengan el mismo número de tokens que puntúan.
\end{observacion}

\section{La plantilla de chat}

Un diálogo es una lista de mensajes $(r_1,c_1)$, \dots, $(r_J,c_J)$, donde $r_j\in\{\text{sistema},\text{usuario},\text{asistente}\}$ es el rol y $c_j$ el contenido. Para tratarlo con un modelo autorregresivo hay que convertirlo en una única secuencia.

\begin{definicion}[Plantilla de chat]\label{def:plantilla}
Sea $\V$ un vocabulario ampliado con un token especial por rol, $\langle r\rangle$, y un token de fin de turno, $\langle\mathrm{fin}\rangle$, que no aparecen en el texto ordinario. La \emph{plantilla de chat} asigna al diálogo la secuencia
\[
\langle r_1\rangle\,\tau(c_1)\,\langle\mathrm{fin}\rangle\ \langle r_2\rangle\,\tau(c_2)\,\langle\mathrm{fin}\rangle\ \cdots\ \langle r_J\rangle\,\tau(c_J)\,\langle\mathrm{fin}\rangle,
\]
donde $\tau$ es el tokenizador. En el ajuste supervisado, la máscara vale $1$ en los tokens de $\tau(c_j)$ y en el $\langle\mathrm{fin}\rangle$ que los cierra si $r_j=\text{asistente}$, y $0$ en el resto.
\end{definicion}

\addtocounter{figure}{1}%

En la inferencia, la aplicación construye la plantilla con todo el historial, añade $\langle\text{asistente}\rangle$ y deja que el modelo genere hasta emitir $\langle\mathrm{fin}\rangle$. Incluir el token de fin en la pérdida es imprescindible para que el modelo aprenda a terminar su turno: en la práctica del módulo, el modelo base no lo emite nunca y, tras 600 pasos de SFT, lo emite en el $99{,}4\,\%$ de los ejemplos de entrenamiento y en el $100\,\%$ de los de prueba.

\begin{observacion}
Los tokens especiales se añaden al vocabulario antes del ajuste (módulo~12) y sus vectores de \emph{embedding} se aprenden durante él. Por eso la plantilla forma parte del modelo: usar un modelo ajustado con una plantilla distinta de la suya degrada su comportamiento. Cada laboratorio usa su propio formato interno y no siempre lo publica; lo que documentan públicamente las API de Claude, Gemini y Grok es que los mensajes se envían como una lista de pares (rol, contenido).
\end{observacion}

\section{¿Cuántos datos hacen falta? La alineación superficial}

El preentrenamiento consume billones de tokens; el ajuste supervisado publicado usa entre miles y cientos de miles de diálogos. Zhou \emph{et al.}~\cite{zhou-lima} ajustaron un modelo LLaMa de 65\,000 millones de parámetros con solo 1000 pares instrucción--respuesta seleccionados con cuidado (LIMA); en una evaluación humana controlada, sus respuestas empataron con las de GPT-4 o fueron preferidas a ellas en el $43\,\%$ de los casos. De ahí la \emph{hipótesis de la alineación superficial}: casi todo el conocimiento del modelo se adquiere en el preentrenamiento, y el ajuste enseña sobre todo qué formato y qué estilo usar al interactuar con un usuario. En la misma línea, Dettmers \emph{et al.}~\cite[§1]{dettmers} observaron que un conjunto de 9000 diálogos de alta calidad (OASST1) daba mejores asistentes que una submuestra de 450\,000 ejemplos de FLAN~v2.

\begin{observacion}
Estas conclusiones proceden de evaluaciones concretas de calidad conversacional y no implican que los datos de SFT sean irrelevantes para todas las capacidades: tareas como el razonamiento matemático o la programación se benefician de muchos más ejemplos. La lección robusta es que la calidad y la diversidad de los datos pesan más que su cantidad.
\end{observacion}


% =====================================================================
\chapter{Olvido catastrófico}
\label{tema:3}
% =====================================================================

Al ajustar todos los pesos con una tarea estrecha, el modelo pierde capacidades que tenía. Este tema define el fenómeno, lo analiza en un modelo cuadrático que puede resolverse a mano y deduce la penalización de la consolidación elástica de pesos (EWC) como aproximación bayesiana. El fenómeno se describió en redes conexionistas en el trabajo de McCloskey y Cohen~\cite{mccloskey}; EWC procede de Kirkpatrick \emph{et al.}~\cite{kirkpatrick}; la aproximación de Laplace sigue el §4.4 de Bishop~\cite{bishop}, y el análisis cuadrático de la penalización, el §7.1.1 de Goodfellow, Bengio y Courville~\cite{gbc}.

\section{El fenómeno}

\begin{definicion}[Olvido]\label{def:olvido}
Sean $\mathcal{L}_A$ la pérdida de una tarea antigua, $\theta_0$ unos parámetros que la resuelven y $\theta_B$ los parámetros obtenidos al ajustar $\theta_0$ con una tarea nueva de pérdida $\mathcal{L}_B$. El \emph{olvido} es el aumento $\mathcal{L}_A(\theta_B)-\mathcal{L}_A(\theta_0)$. Se habla de \emph{olvido catastrófico} cuando ese aumento es comparable al que tendría un modelo sin entrenar en la tarea~$A$.
\end{definicion}

McCloskey y Cohen~\cite{mccloskey} mostraron en 1989 que una red entrenada secuencialmente en dos tareas pierde casi por completo la primera al aprender la segunda, a diferencia de lo que ocurre en el aprendizaje humano. La causa es que los mismos parámetros codifican ambas tareas y el gradiente de $\mathcal{L}_B$ no tiene en cuenta cuánto importa cada parámetro para~$A$.

\addtocounter{definicion}{1}%

\section{Un modelo cuadrático del compromiso}

\begin{proposicion}[Muelle hacia el modelo preentrenado]\label{prop:muelle}
Sean $a,\theta_0\in\R$ y $\lambda\ge 0$, y sea
\[
\mathcal{L}(\theta)=(\theta-a)^{2}+\lambda(\theta-\theta_0)^{2}.
\]
El único minimizador es $\theta^{\ast}=(a+\lambda\theta_0)/(1+\lambda)$. En él, el olvido medido como $(\theta^{\ast}-\theta_0)^{2}$ y el error en la tarea nueva $(\theta^{\ast}-a)^{2}$ valen
\[
(\theta^{\ast}-\theta_0)^{2}=\frac{(a-\theta_0)^{2}}{(1+\lambda)^{2}},\qquad (\theta^{\ast}-a)^{2}=\frac{\lambda^{2}(a-\theta_0)^{2}}{(1+\lambda)^{2}}.
\]
\end{proposicion}

\begin{corolario}\label{cor:frontera}
Las distancias $\abs{\theta^{\ast}-\theta_0}$ y $\abs{\theta^{\ast}-a}$ suman $\abs{a-\theta_0}$ para todo $\lambda\ge 0$: $\theta^{\ast}$ recorre el segmento entre $a$ y $\theta_0$, desde $\theta^{\ast}=a$ ($\lambda=0$, olvido máximo) hasta $\theta^{\ast}\to\theta_0$ ($\lambda\to\infty$, ningún aprendizaje).
\end{corolario}

\addtocounter{definicion}{1}%

La penalización $\lambda(\theta-\theta_0)^{2}$ es la energía potencial de un muelle de constante $2\lambda$ anclado en $\theta_0$, y $\theta^{\ast}$ es el punto de equilibrio entre dos fuerzas lineales. El problema real tiene millones de parámetros, y no todos importan igual para la tarea antigua: conviene un muelle más rígido en los parámetros importantes. Esa es la idea de EWC.

\section{Consolidación elástica de pesos}

\subsection{Justificación bayesiana}

Sean $\mathcal{D}_A$ y $\mathcal{D}_B$ los datos de las dos tareas, condicionalmente independientes dado $\theta$. Por el teorema de Bayes,
\begin{equation}\label{eq:bayes-secuencial}
\ln p(\theta\mid\mathcal{D}_A,\mathcal{D}_B)=\ln p(\mathcal{D}_B\mid\theta)+\ln p(\theta\mid\mathcal{D}_A)-\ln p(\mathcal{D}_B\mid\mathcal{D}_A).
\end{equation}
El último término no depende de $\theta$. Toda la información sobre la tarea~$A$ está en la distribución a posteriori $p(\theta\mid\mathcal{D}_A)$, que no puede calcularse exactamente. La \emph{aproximación de Laplace}~\cite[§4.4]{bishop} la sustituye por una gaussiana centrada en su moda.

\begin{definicion}[Aproximación de Laplace]\label{def:laplace}
Sea $\theta_A$ un máximo de $g(\theta)=\ln p(\theta\mid\mathcal{D}_A)$ y sea $H_A=-\nabla^{2}g(\theta_A)$, definida positiva. La aproximación de Laplace es $p(\theta\mid\mathcal{D}_A)\approx\mathcal{N}(\theta;\theta_A,H_A^{-1})$, que se obtiene del desarrollo de Taylor de segundo orden $g(\theta)\approx g(\theta_A)-\frac12(\theta-\theta_A)\T H_A(\theta-\theta_A)$, en el que el término lineal se anula porque $\nabla g(\theta_A)=0$.
\end{definicion}

Sustituyendo en~\eqref{eq:bayes-secuencial}, maximizar la distribución a posteriori conjunta equivale aproximadamente a minimizar
\[
-\ln p(\mathcal{D}_B\mid\theta)+\tfrac12(\theta-\theta_A)\T H_A(\theta-\theta_A).
\]
Queda por estimar $H_A$, una matriz de $D\times D$ con $D$ del orden de millones.

\subsection{Información de Fisher}

\begin{definicion}[Información de Fisher]\label{def:fisher}
Sea $p_\theta$ una familia de distribuciones derivable en $\theta$. Su \emph{matriz de información de Fisher} es
\[
\Fis(\theta)=\E_{x\sim p_\theta}\bigl[\nabla_\theta\ln p_\theta(x)\,\nabla_\theta\ln p_\theta(x)\T\bigr].
\]
El vector $\nabla_\theta\ln p_\theta(x)$ se llama \emph{función de puntuación} (\emph{score}).
\end{definicion}

$\Fis(\theta)$ es semidefinida positiva por ser la esperanza de matrices de la forma $vv\T$. Solo requiere gradientes de primer orden.

\begin{proposicion}[Fisher como hessiano esperado]\label{prop:fisher-hess}
Supóngase que se pueden intercambiar derivada e integral. Entonces $\E_{x\sim p_\theta}[\nabla_\theta\ln p_\theta(x)]=0$ y
\[
\Fis(\theta)=-\E_{x\sim p_\theta}\bigl[\nabla^{2}_\theta\ln p_\theta(x)\bigr].
\]
\end{proposicion}

Si $\mathcal{D}_A$ consta de $n$ ejemplos y la distribución a priori es plana, $H_A=-\sum_{i}\nabla^{2}\ln p_\theta(x_i)\big|_{\theta_A}$, y por la ley de los grandes números y la \cref{prop:fisher-hess} se aproxima por $n\,\Fis(\theta_A)$ cuando el modelo describe bien los datos. Kirkpatrick \emph{et al.} se quedan además con la diagonal, $F_i=\Fis_{ii}(\theta_A)$, que se estima promediando los cuadrados de las componentes del gradiente.

\begin{definicion}[EWC]\label{def:ewc}
Sean $\theta_A$ los parámetros tras la tarea~$A$ y $F_i$ la diagonal de su matriz de Fisher. La \emph{consolidación elástica de pesos} minimiza
\[
\mathcal{L}_{\mathrm{EWC}}(\theta)=\mathcal{L}_B(\theta)+\sum_{i=1}^{D}\frac{\lambda}{2}F_i\,(\theta_i-\theta_{A,i})^{2},
\]
donde $\lambda>0$ regula la importancia de la tarea antigua~\cite{kirkpatrick}.
\end{definicion}

Cada parámetro tiene su propio muelle, de constante $\lambda F_i$: los parámetros a los que la verosimilitud de la tarea~$A$ es sensible (Fisher grande) quedan rígidos y los demás pueden moverse libremente. La \cref{prop:muelle} es el caso de un solo parámetro en el que el coeficiente $\frac{\lambda}{2}F$ de EWC hace el papel del $\lambda$ de la proposición.

\section{Otros remedios}

\begin{definicion}[Repetición]\label{def:replay}
La \emph{repetición} (\emph{replay}) mezcla en cada paso ejemplos de la tarea antigua: con $\beta\in[0,1)$, se minimiza $(1-\beta)\mathcal{L}_B(\theta)+\beta\,\mathcal{L}_A(\theta)$.
\end{definicion}

Otros remedios habituales son una tasa de aprendizaje pequeña y pocas épocas, las técnicas de ULMFiT (\cref{tema:1}) y modificar pocos parámetros, como hacen los métodos del \cref{tema:4} y LoRA (\cref{tema:5}). El módulo~21 penalizará la divergencia de Kullback--Leibler entre el modelo ajustado y uno de referencia. Esa penalización actúa en el espacio de las distribuciones de salida, pero localmente coincide con un muelle en el espacio de parámetros cuya rigidez es la información de Fisher.

\begin{proposicion}[Desarrollo local de la divergencia KL]\label{prop:kl-fisher}
Bajo las hipótesis de la \cref{prop:fisher-hess}, y si $\ln p_\theta$ es tres veces derivable,
\[
D_{\mathrm{KL}}\bigl(p_{\theta_0}\,\big\|\,p_{\theta_0+\delta}\bigr)=\tfrac12\,\delta\T\Fis(\theta_0)\,\delta+O(\norm{\delta}^{3}).
\]
\end{proposicion}


% =====================================================================
\chapter{Ajuste eficiente en parámetros}
\label{tema:4}
% =====================================================================

El ajuste completo de un modelo de $D$ parámetros produce, para cada tarea, una copia de $D$ parámetros y exige guardar durante el entrenamiento el estado del optimizador de todos ellos. Los métodos de \emph{ajuste eficiente en parámetros} (PEFT, del inglés \emph{parameter-efficient fine-tuning}) congelan el modelo preentrenado y entrenan un número pequeño de parámetros nuevos. Este tema presenta los adaptadores, el ajuste de \emph{prompt} y de prefijo, y la noción de dimensión intrínseca que explica por qué bastan tan pocos parámetros. Se sigue el §8.2 de Jurafsky y Martin~\cite{jm3}, la clase~9 de CS224N~\cite{cs224n} y los artículos originales~\cite{houlsby,lester,li-liang,li-intrinsic,aghajanyan}; LoRA, el método más usado, ocupa el \cref{tema:5}.

\section{Planteamiento}

\begin{definicion}[Ajuste eficiente en parámetros]\label{def:peft}
Sea $\theta_0\in\R^{D}$ un modelo preentrenado. Un método PEFT define una familia de modelos $f_{\theta_0,\varphi}$ con parámetros nuevos $\varphi\in\R^{P}$, $P\ll D$, tal que existe $\varphi_0$ con $f_{\theta_0,\varphi_0}=f_{\theta_0}$, y entrena solo $\varphi$, partiendo de $\varphi_0$ o de un valor próximo, con $\theta_0$ congelado.
\end{definicion}

Los métodos se agrupan en tres familias: \emph{aditivos}, que añaden módulos o vectores nuevos (adaptadores, ajuste de \emph{prompt} y de prefijo); \emph{selectivos}, que entrenan un subconjunto de los parámetros existentes (la sonda lineal del \cref{tema:1} es el caso extremo), y de \emph{reparametrización}, que escriben la actualización $\Delta\theta$ como función de pocos parámetros (subespacios aleatorios, \cref{sec:intrinseca}, y LoRA, \cref{tema:5}).

\begin{proposicion}[Memoria del entrenamiento]\label{prop:memoria-peft}
Con precisión mixta y Adam, el ajuste completo requiere unos $16D$ bytes para pesos, gradientes y estado del optimizador (módulo~17). Con un método PEFT y el modelo base en bf16, la memoria es de unos $2D+16P$ bytes, sin contar las activaciones.
\end{proposicion}

Con $D=7\cdot 10^{9}$ y $P=2\cdot 10^{7}$, la memoria baja de $112$~GB a $14+0{,}32=14{,}32$~GB. Además, cada tarea se almacena con $P$ parámetros en lugar de $D$, y un mismo modelo base puede servir muchas tareas cambiando solo $\varphi$.

\section{Adaptadores}

\begin{definicion}[Adaptador]\label{def:adaptador}
Sea $h\in\R^{d}$ la salida de una subcapa de un Transformer. Un \emph{adaptador} con dimensión de cuello $m\ll d$ calcula
\[
a(h)=h+W_{\uparrow}\,g\bigl(W_{\downarrow}h+b_{\downarrow}\bigr)+b_{\uparrow},\qquad W_{\downarrow}\in\R^{m\times d},\ W_{\uparrow}\in\R^{d\times m},
\]
con $g$ una no linealidad. Houlsby \emph{et al.}~\cite{houlsby} insertan dos adaptadores por capa, tras la subcapa de atención y tras la red de alimentación hacia delante, y entrenan además los parámetros de normalización de capa.
\end{definicion}

\addtocounter{figure}{1}%

\begin{proposicion}\label{prop:adaptador}
Un adaptador tiene $2md+m+d$ parámetros. Si se inicializa con $W_{\uparrow}=0$ y $b_{\uparrow}=0$, entonces $a(h)=h$ para todo $h$, de modo que el modelo con adaptadores coincide al principio con el preentrenado.
\end{proposicion}

Houlsby \emph{et al.} usan una inicialización \emph{casi} nula, de modo que el adaptador es casi la identidad, y obtienen en GLUE resultados a menos de un $0{,}4\,\%$ de los del ajuste completo añadiendo un $3{,}6\,\%$ de parámetros por tarea~\cite{houlsby}.

\begin{observacion}[Latencia]
Los adaptadores añaden capas que se ejecutan en serie y no pueden fusionarse con los pesos originales, porque contienen una no linealidad. Hu \emph{et al.}~\cite[§3]{hu-lora} midieron aumentos de la latencia de inferencia superiores al $30\,\%$ con secuencias cortas y lotes pequeños. LoRA evita este problema (\cref{tema:5}).
\end{observacion}

\section{Ajuste de \emph{prompt} y de prefijo}

\begin{definicion}[Ajuste de \emph{prompt}]\label{def:prompt-tuning}
Sea $E\in\R^{T\times d}$ la matriz de \emph{embeddings} de una entrada. El \emph{ajuste de prompt} antepone $p$ vectores entrenables $P_e\in\R^{p\times d}$ y alimenta el modelo congelado con $[P_e;E]\in\R^{(p+T)\times d}$. Tiene $pd$ parámetros~\cite{lester}.
\end{definicion}

Los vectores de $P_e$ no tienen por qué corresponder a ningún token del vocabulario: forman un \emph{prompt continuo} (\emph{soft prompt}). Lester \emph{et al.}~\cite{lester} observaron que la diferencia con el ajuste completo disminuye con la escala y prácticamente desaparece en modelos de unos $10^{10}$ parámetros.

\begin{definicion}[Ajuste de prefijo]\label{def:prefix}
El \emph{ajuste de prefijo} antepone, en cada una de las $L$ capas, $p$ claves y $p$ valores entrenables, $P_K^{l},P_V^{l}\in\R^{p\times d}$, a las claves y valores calculados a partir de la entrada. Tiene $2Lpd$ parámetros. Para estabilizar el entrenamiento, Li y Liang~\cite{li-liang} reparametrizan los prefijos mediante una pequeña red que se descarta al terminar; con ello entrenan en torno al $0{,}1\,\%$ de los parámetros.
\end{definicion}

El efecto de un prefijo sobre la atención admite una descripción exacta.

\begin{teorema}[El prefijo interpola la atención]\label{teo:prefijo}
Sean $q\in\R^{d}$ una consulta, $K,V\in\R^{T\times d}$ las claves y valores de la entrada y $P_K,P_V\in\R^{p\times d}$ los del prefijo. Escríbase $\mathrm{At}(q,K,V)=\softmax(Kq/\sqrt{d})\T V$. Entonces
\[
\mathrm{At}\bigl(q,[P_K;K],[P_V;V]\bigr)=\bigl(1-\mu(q)\bigr)\,\mathrm{At}(q,K,V)+\mu(q)\,\mathrm{At}(q,P_K,P_V),
\]
con
\[
\mu(q)=\frac{S_P}{S_P+S_K},\qquad S_P=\sum_{j=1}^{p}e^{(P_K)_j\cdot q/\sqrt{d}},\quad S_K=\sum_{i=1}^{T}e^{k_i\cdot q/\sqrt{d}}.
\]
\end{teorema}

\begin{corolario}\label{cor:prefijo}
El ajuste de prefijo no altera los pesos relativos de atención entre las posiciones de la entrada: solo mezcla la salida original con un vector que depende de la consulta, con una proporción $\mu(q)\in(0,1)$.
\end{corolario}

Esta lectura, que relaciona el ajuste de prefijo con los adaptadores, se desarrolla en He \emph{et al.}~\cite{he-unified}.

\section{Dimensión intrínseca}\label{sec:intrinseca}

¿Por qué basta con ajustar una fracción mínima de los parámetros? Li \emph{et al.}~\cite{li-intrinsic} propusieron medir la dimensión efectiva de un problema de optimización restringiendo el entrenamiento a un subespacio aleatorio.

\begin{definicion}[Subespacio aleatorio y dimensión intrínseca]\label{def:intrinseca}
Sean $\theta_0\in\R^{D}$ y $P\in\R^{D\times d'}$ una matriz aleatoria fija. Se entrena solo $z\in\R^{d'}$, inicializado en $0$, con la reparametrización $\theta=\theta_0+Pz$. La \emph{dimensión intrínseca} $d_{90}$ de una tarea es el menor $d'$ con el que se alcanza el $90\,\%$ del rendimiento del ajuste completo~\cite[§3]{aghajanyan}.
\end{definicion}

\begin{proposicion}[Gradiente en el subespacio]\label{prop:grad-subespacio}
Si $\mathcal{L}$ es derivable, $\nabla_z\mathcal{L}(\theta_0+Pz)=P\T\,\nabla_\theta\mathcal{L}(\theta)\big|_{\theta=\theta_0+Pz}$.
\end{proposicion}

Aghajanyan \emph{et al.}~\cite{aghajanyan} estimaron la dimensión intrínseca del ajuste de modelos de lenguaje preentrenados, usando para $P$ la transformada Fastfood, que evita almacenar una matriz densa de $D\times d'$. Sus tres resultados principales son:
\begin{enumerate}
  \item con RoBERTa-Large, unos 200 parámetros bastan para alcanzar el $90\,\%$ del rendimiento del ajuste completo en la tarea MRPC~\cite[§4.2]{aghajanyan};
  \item la dimensión intrínseca disminuye a lo largo del preentrenamiento y es menor en los modelos con más parámetros~\cite[§5.1 y~§5.2]{aghajanyan};
  \item una cota de generalización basada en compresión depende de la dimensión intrínseca y no del número total de parámetros: con $n$ ejemplos, el error de generalización es $O(\sqrt{d_{90}/n})$, para modelos entrenados en el subespacio~\cite[§5.3]{aghajanyan}.
\end{enumerate}

\begin{observacion}
El segundo resultado explica en parte la paradoja de que los modelos mayores se ajusten con menos datos: el preentrenamiento comprime el problema de adaptación. Hu \emph{et al.} formularon a partir de aquí la hipótesis que motiva LoRA~\cite[§1 y~§4.1]{hu-lora}: si el ajuste vive en un subespacio de dimensión pequeña, también la actualización de cada matriz de pesos debe tener un \emph{rango intrínseco} bajo. LoRA no fija las direcciones al azar: aprende la mejor base de rango $r$.
\end{observacion}


% =====================================================================
\chapter{LoRA: adaptación de bajo rango}
\label{tema:5}
% =====================================================================

LoRA (\emph{low-rank adaptation}) congela cada matriz preentrenada y aprende solo una corrección de rango bajo, escrita como producto de dos matrices delgadas. Es hoy el método PEFT más usado porque no añade latencia de inferencia y su análisis se reduce al álgebra lineal del módulo~3. Este tema sigue el artículo original de Hu \emph{et al.}~\cite[§4 y~§7]{hu-lora}, el teorema de Eckart y Young~\cite{eckart-young} en la forma del módulo~3 y de Deisenroth, Faisal y Ong~\cite[§4.6]{mml}, y la clase invitada sobre LoRA del curso CS224N~\cite{cs224n}.

\section{Definición y número de parámetros}

\begin{definicion}[LoRA]\label{def:lora}
Sea $W\in\R^{d\times k}$ una matriz preentrenada de una capa lineal $h=Wx$, con $x\in\R^{k}$. Sean $r\ll\min(d,k)$ el \emph{rango}, $\alpha>0$ un hiperparámetro de escala y $s=\alpha/r$. LoRA sustituye la capa por
\[
h=Wx+s\,BAx,\qquad A\in\R^{r\times k},\ B\in\R^{d\times r},
\]
con $W$ congelada y $A$, $B$ entrenables. Se inicializa $A$ con entradas gaussianas y $B=0$~\cite[§4.1]{hu-lora}. La matriz efectiva es $W'=W+s\,BA$.
\end{definicion}

\addtocounter{figure}{1}%

\begin{proposicion}[Parámetros y coste]\label{prop:lora-param}
LoRA entrena $r(d+k)$ parámetros por matriz, una fracción $r(d+k)/(dk)=r(1/d+1/k)$ de los $dk$ del ajuste completo. Calcular $s\,B(Ax)$ cuesta $2r(d+k)$ operaciones más un escalado, frente a las $2dk$ de $Wx$, y nunca requiere formar $BA$.
\end{proposicion}

\begin{table}[h]
\centering\small
\caption{Parámetros entrenables de LoRA en una matriz.}\label{tab:lora}
\begin{tabular}{@{}lrrrr@{}}
\toprule
Matriz & $r$ & Completo $dk$ & LoRA $r(d+k)$ & Fracción \\
\midrule
$64\times 64$ & 4 & $4096$ & $512$ & $12{,}5\,\%$ \\
$768\times 768$ & 4 & $589\,824$ & $6144$ & $1{,}04\,\%$ \\
$4096\times 4096$ & 8 & $16\,777\,216$ & $65\,536$ & $0{,}39\,\%$ \\
$4096\times 4096$ & 16 & $16\,777\,216$ & $131\,072$ & $0{,}78\,\%$ \\
\bottomrule
\end{tabular}
\end{table}

\addtocounter{definicion}{1}%

En el mini-GPT de la práctica ($d=64$, red de alimentación de dimensión 256, 2 bloques), LoRA con $r=4$ en las cuatro matrices de atención ($64\times 64$) y las dos de la red de alimentación ($64\times 256$ y $256\times 64$) de cada bloque entrena $2\cdot(4\cdot 512+2\cdot 1280)=9216$ parámetros, el $8{,}57\,\%$ de los $107\,520$ del modelo. En GPT-3 de 175\,000 millones de parámetros, Hu \emph{et al.} redujeron los parámetros entrenables unas $10\,000$ veces y la memoria de GPU unas 3 veces respecto al ajuste completo con Adam; con $r=4$ en $W_Q$ y $W_V$, el punto de control pasó de unos 350~GB a unos 35~MB~\cite[§4.2]{hu-lora}.

\begin{proposicion}[Fusión]\label{prop:fusion}
Tras el entrenamiento, sustituir $W$ por $W'=W+s\,BA$ produce la misma función que la capa LoRA y el mismo coste de inferencia que la capa original.
\end{proposicion}

Así, LoRA no añade latencia de inferencia, a diferencia de los adaptadores del \cref{tema:4}. Sin fusionar, en cambio, un mismo modelo base puede servir muchos adaptadores pequeños e intercambiarlos en cada petición.

\section{Rango de la corrección}

\begin{proposicion}[Cota de rango]\label{prop:rango}
Para todo $A\in\R^{r\times k}$ y $B\in\R^{d\times r}$, $\rango(BA)\le r$.
\end{proposicion}

\begin{proposicion}[Recíproco]\label{prop:factorizacion}
Toda matriz $M\in\R^{d\times k}$ con $\rango M\le r$ puede escribirse como $M=BA$ con $B\in\R^{d\times r}$ y $A\in\R^{r\times k}$.
\end{proposicion}

\begin{corolario}\label{cor:representables}
El conjunto de correcciones que LoRA de rango $r$ puede representar es exactamente $\{\Delta W\in\R^{d\times k}:\rango\Delta W\le r\}$. Si $r<\min(d,k)$, ninguna corrección de rango completo es representable.
\end{corolario}

\section{Gradientes e inicialización}

\begin{proposicion}[Gradientes de LoRA]\label{prop:grad-lora}
Sea $h=Wx+s\,BAx$ y sea $G=\partial\mathcal{L}/\partial h\in\R^{d}$ el gradiente que llega a la salida. Entonces
\[
\frac{\partial\mathcal{L}}{\partial B}=s\,G\,(Ax)\T\in\R^{d\times r},\qquad
\frac{\partial\mathcal{L}}{\partial A}=s\,B\T G\,x\T\in\R^{r\times k}.
\]
Con un lote, se suman las contribuciones de cada ejemplo.
\end{proposicion}

\begin{corolario}[Inicialización]\label{cor:init}
Con $B=0$ y $A$ aleatoria:
\begin{enumerate}
  \item $BA=0$ y el modelo ajustado empieza siendo exactamente el modelo base;
  \item en el primer paso $\partial\mathcal{L}/\partial A=0$, pero $\partial\mathcal{L}/\partial B=s\,G(Ax)\T\ne 0$ si $G\ne 0$ y $Ax\ne 0$, de modo que $B$ deja de ser nula y, desde el segundo paso, también $A$ recibe gradiente.
\end{enumerate}
En cambio, $(A,B)=(0,0)$ es un punto fijo del descenso de gradiente (\cref{eje:5-init}).
\end{corolario}

La elección simétrica, $A=0$ con $B$ aleatoria, también empieza en el modelo base y también permite aprender; el convenio del artículo original es $B=0$.

\section{La escala \texorpdfstring{$\alpha/r$}{alpha/r}}

\begin{proposicion}[Adam es invariante a la escala del gradiente]\label{prop:adam-escala}
Sea el optimizador Adam con $\varepsilon=0$ (módulo~7). Si se multiplican todos los gradientes por una constante $c>0$, la sucesión de actualizaciones de los parámetros no cambia.
\end{proposicion}

Por la \cref{prop:grad-lora}, el gradiente de $B$ es proporcional a $s$. Por la \cref{prop:adam-escala}, la actualización $\Delta B$ del primer paso no depende de $s$, y la variación de la función, $s\,\Delta B\,A\,x$, es proporcional a $s$. El factor $s=\alpha/r$ actúa, por tanto, como un multiplicador de la tasa de aprendizaje. Hu \emph{et al.} observan que ajustar $\alpha$ con Adam equivale aproximadamente a ajustar la tasa, fijan $\alpha$ al primer valor de $r$ que prueban y no lo modifican; dividir por $r$ reduce la necesidad de reajustar hiperparámetros al cambiar el rango~\cite[§4.1]{hu-lora}. En la práctica se fija $\alpha$ (por ejemplo, $\alpha=16$ o $\alpha=2r$) y se ajusta $r$.

\section{Por qué funciona: relación con la SVD}

LoRA apuesta a que la actualización $\Delta W$ que necesita el ajuste fino tiene un espectro que decae deprisa. El teorema de Eckart--Young (módulo~3, tema~5) cuantifica lo que se pierde al restringir el rango.

\begin{teorema}[Eckart--Young, norma de Frobenius]\label{teo:ey}
Sea $\Delta W=\sum_i\sigma_iu_iv_i\T$ la SVD, con $\sigma_1\ge\sigma_2\ge\cdots$. Para todo $r$, la matriz $\Delta W_r=\sum_{i\le r}\sigma_iu_iv_i\T$ minimiza $\norm{\Delta W-M}_F$ entre las matrices de rango a lo sumo~$r$, y
\[
\norm{\Delta W-\Delta W_r}_F=\Bigl(\sum_{i>r}\sigma_i^{2}\Bigr)^{1/2}.
\]
\end{teorema}

La demostración está en el módulo~3 y en~\cite{eckart-young}.

\begin{corolario}[Error mínimo de LoRA]\label{cor:error-lora}
Sea $\Delta W^{\star}$ la actualización ideal de una matriz. Entonces
\[
\min_{A,B}\,\norm{s\,BA-\Delta W^{\star}}_F=\Bigl(\sum_{i>r}\sigma_i^{2}(\Delta W^{\star})\Bigr)^{1/2}.
\]
\end{corolario}

\addtocounter{definicion}{1}%

Hay dos diferencias con truncar una SVD. LoRA no calcula ninguna SVD: el descenso de gradiente encuentra $B$ y $A$ directamente. Y LoRA aproxima la \emph{actualización}, no $W$, que suele tener rango completo y no se puede comprimir así sin perder mucho.

\begin{observacion}[Evidencia empírica]
Hu \emph{et al.} adaptan principalmente $W_Q$ y $W_V$ y encuentran que, en algunas tareas, un rango tan pequeño como $r=1$ basta~\cite[§7.1 y~§7.2]{hu-lora}; las direcciones singulares principales aprendidas con $r=8$ y con $r=64$ se solapan mucho, y el resto parecen ruido. Dettmers \emph{et al.} observaron, en cambio, que para igualar el ajuste completo en modelos grandes hace falta aplicar LoRA a todas las capas lineales~\cite[§4]{dettmers}. En la práctica del módulo, la actualización del ajuste completo en $W_Q$ tiene rango 64 y sus 4 primeros valores singulares solo concentran el $41{,}5\,\%$ de la energía: un modelo tan pequeño, ajustado de forma tan agresiva, no cumple bien la hipótesis. Aun así, LoRA acierta más en la prueba ($0{,}450$ frente a $0{,}300$, con tasas de aprendizaje distintas); una explicación plausible es que restringir el rango regulariza.
\end{observacion}


% =====================================================================
\chapter{QLoRA y elección de estrategia}
\label{tema:6}
% =====================================================================

Con LoRA, la matriz $W$ no recibe actualizaciones: solo se lee. Puede guardarse, por tanto, en muy baja precisión. QLoRA combina la cuantización del modelo base a 4 bits con adaptadores LoRA entrenados en 16 bits y permite ajustar modelos de decenas de miles de millones de parámetros en una sola GPU. Este tema estudia sus tres ingredientes (el formato NF4, la doble cuantización y los optimizadores paginados), hace la contabilidad de memoria de cada configuración y termina con los criterios para elegir entre ajuste fino, instrucciones en el contexto y recuperación de documentos. Se sigue el artículo de Dettmers \emph{et al.}~\cite[§3 y~§4]{dettmers}, la contabilidad de parámetros de la tarea~1 de CS336~\cite[§3.5]{cs336a1} y el capítulo~8 de Jurafsky y Martin~\cite{jm3}; la cuantización por bloques se estudió en el módulo~19.

\section{Cuantización por bloques}

\begin{definicion}[Cuantización por bloques con un libro de códigos]\label{def:cuant-bloques}
Sea $\mathcal{Q}=\{q_0<q_1<\dots<q_{2^{b}-1}\}\subset[-1,1]$ un \emph{libro de códigos} de $2^{b}$ niveles. Los pesos se dividen en bloques de $B$ elementos. Para cada bloque $w\in\R^{B}$ se guarda la \emph{constante de escala} $c=\max_j\abs{w_j}$ y, para cada $j$, el índice del nivel de $\mathcal{Q}$ más próximo a $w_j/c$. La \emph{descuantización} devuelve $\hat w_j=c\,q_{i_j}$.
\end{definicion}

\begin{proposicion}[Bits por parámetro]\label{prop:bits}
Si cada índice ocupa $b$ bits y cada constante se guarda en coma flotante de 32 bits, la cuantización por bloques ocupa $b+32/B$ bits por parámetro.
\end{proposicion}

Con $b=4$ y $B=64$, el coste es $4+0{,}5=4{,}5$ bits por parámetro: la constante de escala añade un $12{,}5\,\%$ al tamaño de los índices.

\section{El formato NF4}

La elección del libro de códigos importa. Los pesos de una red preentrenada siguen aproximadamente una distribución normal de media nula, de modo que los valores próximos a cero son mucho más frecuentes que los extremos. Un libro de códigos uniforme desperdicia niveles en zonas poco pobladas.

\begin{proposicion}[Cuantización por cuantiles]\label{prop:cuantiles}
Sea $X\sim\mathcal{N}(0,1)$ y sean $\beta_i=\Phi^{-1}(i/2^{b})$, $i=0,\dots,2^{b}$, con $\beta_0=-\infty$ y $\beta_{2^{b}}=+\infty$. Si se asigna el código $i$ a los valores de $[\beta_i,\beta_{i+1})$, cada código tiene probabilidad $2^{-b}$, y la entropía del código es $b$ bits, la máxima posible con $2^{b}$ símbolos.
\end{proposicion}

Así, cada uno de los $2^{b}$ códigos se usa con la misma frecuencia y ninguno se desperdicia. El \emph{NormalFloat} de 4 bits (NF4) de Dettmers \emph{et al.} aplica esta idea con $b=4$: estima los niveles a partir de cuantiles de la normal estándar, los normaliza al intervalo $[-1,1]$ y, para que el cero sea representable exactamente, construye por separado $2^{b-1}$ niveles negativos y $2^{b-1}+1$ positivos y elimina el cero repetido~\cite[§3]{dettmers}. El tamaño de bloque es $B=64$.

\addtocounter{definicion}{1}%

\section{Doble cuantización}

\begin{proposicion}[Coste de las constantes con doble cuantización]\label{prop:doble}
Si las constantes de escala de primer nivel, una por bloque de $B_1$ pesos, se cuantizan a su vez a 8 bits en bloques de $B_2$ constantes, cada una con una constante de segundo nivel de 32 bits, el coste de las constantes es $8/B_1+32/(B_1B_2)$ bits por parámetro.
\end{proposicion}

Con $B_1=64$ y $B_2=256$, los valores de QLoRA~\cite[§3]{dettmers}, el coste pasa de $32/64=0{,}5$ a $8/64+32/16\,384=0{,}125+0{,}00195=0{,}127$ bits por parámetro: un ahorro de $0{,}373$ bits por parámetro. Antes de cuantizarlas, las constantes se centran restando su media, lo que permite una cuantización simétrica.

\section{QLoRA}

\begin{definicion}[QLoRA]\label{def:qlora}
QLoRA almacena los pesos congelados del modelo base en NF4 con doble cuantización y añade adaptadores LoRA en bf16 a las capas lineales. La pasada hacia delante de una capa es
\[
Y=X\,\mathrm{dq}(c_1,c_2,W_{\mathrm{NF4}})+s\,X L_1L_2,
\]
donde $\mathrm{dq}$ descuantiza primero las constantes y después los pesos, produciendo una matriz en bf16, y $L_1,L_2$ son los factores LoRA~\cite[§3, ecs.~5 y~6]{dettmers}. El tipo de almacenamiento es NF4 y el tipo de cálculo, bf16. Solo se calculan gradientes respecto a los parámetros de LoRA.
\end{definicion}

En esta definición se usa el convenio de filas de~\cite{dettmers}, $Y=XW$; los factores $L_1,L_2$ desempeñan el papel de $A\T$ y $B\T$. Aunque $W$ no recibe gradiente, la retropropagación sí necesita multiplicar por $W$ descuantizada para propagar el gradiente hacia las capas anteriores, cuyos adaptadores también se entrenan.

El tercer ingrediente son los \emph{optimizadores paginados}: el estado del optimizador se guarda en memoria unificada, que el controlador traslada automáticamente a la memoria de la CPU cuando la GPU se queda sin espacio y devuelve cuando hace falta para la actualización. Con estos tres elementos, Dettmers \emph{et al.} ajustaron un modelo de 65\,000 millones de parámetros en una sola GPU de 48~GB, frente a los más de 780~GB que requería el ajuste completo en 16 bits~\cite[§1]{dettmers}. Su mejor modelo, Guanaco~65B, alcanzó el $99{,}3\,\%$ del rendimiento de ChatGPT en la prueba de Vicuna según la evaluación de GPT-4, con un intervalo de confianza del $95\,\%$ de $\pm 4{,}4$ puntos. Los propios autores advierten que no demostraron que QLoRA iguale al ajuste completo en 16 bits a las escalas de 33\,000 y 65\,000 millones de parámetros~\cite[§8]{dettmers}.

\begin{proposicion}[Memoria de pesos y optimizador]\label{prop:mem-qlora}
Sin contar las activaciones, la memoria de pesos y estado del optimizador es de unos $16N$ bytes en el ajuste completo con Adam en precisión mixta y de unos $(b_{\mathrm{ef}}/8)\,N+16N_{\mathrm{LoRA}}$ bytes en QLoRA, donde $b_{\mathrm{ef}}$ son los bits efectivos por parámetro del modelo base.
\end{proposicion}

\begin{table}[h]
\centering\small
\caption{Memoria de pesos y optimizador de un modelo de $7\cdot 10^{9}$ parámetros, sin activaciones, tomando 4 bits por parámetro para la base cuantizada ($1\,\mathrm{GB}=10^{9}$ bytes).}\label{tab:memoria}
\begin{tabular}{@{}lll@{}}
\toprule
Configuración & Cálculo & Memoria \\
\midrule
Pesos en bf16 & $7\cdot 10^{9}\times 2$ B & $14$ GB \\
Ajuste completo con Adam, precisión mixta & $7\cdot 10^{9}\times 16$ B & $112$ GB \\
QLoRA: base en 4 bits & $7\cdot 10^{9}\times 0{,}5$ B & $3{,}5$ GB \\
QLoRA: $2\cdot 10^{7}$ parámetros LoRA & $2\cdot 10^{7}\times 16$ B & $0{,}32$ GB \\
QLoRA en total & $3{,}5+0{,}32$ & $3{,}82$ GB \\
\bottomrule
\end{tabular}
\end{table}

La reducción es de unas $112/3{,}82\approx 29$ veces y permite ajustar un modelo de este tamaño en una GPU de consumo. Con las constantes de escala, $b_{\mathrm{ef}}\approx 4{,}127$ en lugar de 4, lo que añade unos $0{,}11$~GB.

\section{¿Ajustar, dar instrucciones en el contexto o recuperar documentos?}

El ajuste fino no es la única forma de adaptar un modelo. Las otras dos son dar instrucciones y ejemplos en el contexto (\emph{prompting} y aprendizaje con pocos ejemplos) y recuperar documentos relevantes e insertarlos en el contexto (RAG, módulo~22). La regla práctica es que el ajuste cambia \emph{cómo} responde el modelo y el contexto cambia \emph{qué sabe} en cada consulta.

\begin{table}[h]
\centering\small
\caption{Criterios para elegir la forma de adaptación.}\label{tab:eleccion}
\begin{tabularx}{\textwidth}{@{}>{\raggedright\arraybackslash}p{3.3cm}>{\raggedright\arraybackslash}X>{\raggedright\arraybackslash}X>{\raggedright\arraybackslash}X@{}}
\toprule
Criterio & Instrucciones en el contexto & RAG & Ajuste fino \\
\midrule
Coste inicial & ninguno & índice de documentos & datos y entrenamiento \\
Conocimiento nuevo o cambiante & solo si cabe en el contexto & sí, con fuentes citables & mal: caro de actualizar y propenso a inventar \\
Formato, estilo y tono & aceptable & no ayuda & muy bueno \\
Coste por consulta & contextos largos & recuperación y contexto & contextos cortos \\
Riesgo de olvido & ninguno & ninguno & sí (\cref{tema:3}) \\
Acceso a los pesos & no & no & sí, o una API de ajuste \\
\bottomrule
\end{tabularx}
\end{table}

\addtocounter{figure}{1}%

En la práctica, las tres formas se combinan: un modelo ajustado para seguir un formato y citar sus fuentes, alimentado con documentos recuperados. Los proveedores de modelos comerciales ofrecen públicamente, según el caso, ajuste fino a través de sus plataformas; los detalles técnicos de esos servicios no siempre se publican.



\backmatter
% =====================================================================
\begin{thebibliography}{99}
\addcontentsline{toc}{chapter}{Bibliografía}
\markboth{Bibliografía}{}
\small

\bibitem{aghajanyan}
A.~Aghajanyan, L.~Zettlemoyer y S.~Gupta.
\newblock Intrinsic dimensionality explains the effectiveness of language model fine-tuning.
\newblock En \emph{Proceedings of ACL-IJCNLP}, 2021.
\newblock \url{https://arxiv.org/abs/2012.13255}

\bibitem{bai-cai}
Y.~Bai \emph{et al.}
\newblock Constitutional AI: Harmlessness from AI feedback.
\newblock Anthropic, 2022.
\newblock \url{https://arxiv.org/abs/2212.08073}

\bibitem{bishop}
C.~M. Bishop.
\newblock \emph{Pattern Recognition and Machine Learning}.
\newblock Springer, 2006.
\newblock §4.4, «The Laplace Approximation».

\bibitem{cmu711}
G.~Neubig.
\newblock 11-711: Advanced Natural Language Processing.
\newblock Carnegie Mellon University, otoño de 2024.
\newblock \url{https://www.phontron.com/class/anlp-fall2024/}

\bibitem{cs224n}
D.~Yang y Y.~Choi.
\newblock CS224N: Natural Language Processing with Deep Learning.
\newblock Stanford University, invierno de 2026. Clase~8, «Post-training (RLHF, SFT, DPO)»; clase~9, «Efficient Adaptation (Prompting + PEFT)»; clase invitada de J.~Schulman, «Tinker and LoRA Without Regret».
\newblock \url{https://web.stanford.edu/class/cs224n/}

\bibitem{cs336}
T.~Hashimoto y P.~Liang.
\newblock CS336: Language Modeling from Scratch.
\newblock Stanford University, primavera de 2025. Clase~15, «Alignment -- SFT/RLHF».
\newblock \url{https://cs336.stanford.edu/spring2025}

\bibitem{cs336a1}
Stanford CS336.
\newblock Assignment~1 (basics): Building a Transformer LM.
\newblock Primavera de 2026. §3.5, problema \texttt{transformer\_\allowbreak accounting}.
\newblock \url{https://github.com/stanford-cs336/assignment1-basics}

\bibitem{cs336a5}
Stanford CS336.
\newblock Assignment~5 (alignment): Alignment and Reasoning RL.
\newblock Primavera de 2025. §4.2, problemas \texttt{tokenize\_\allowbreak prompt\_\allowbreak and\_\allowbreak output}, \texttt{masked\_\allowbreak normalize} y \texttt{sft\_\allowbreak microbatch\_\allowbreak train\_\allowbreak step}.
\newblock \url{https://github.com/stanford-cs336/assignment5-alignment}

\bibitem{cs336a5s}
Stanford CS336.
\newblock Assignment~5 Supplement: Instruction Tuning and RLHF.
\newblock Primavera de 2026. §4.2, problema \texttt{data\_\allowbreak loading}.
\newblock \url{https://github.com/stanford-cs336/assignment5-alignment}

\bibitem{mml}
M.~P. Deisenroth, A.~A. Faisal y C.~S. Ong.
\newblock \emph{Mathematics for Machine Learning}.
\newblock Cambridge University Press, 2020. Versión en línea de 2024.
\newblock Cap.~4 «Matrix Decompositions», §4.6 y ejercicio~4.10.
\newblock \url{https://mml-book.github.io/}

\bibitem{dettmers}
T.~Dettmers, A.~Pagnoni, A.~Holtzman y L.~Zettlemoyer.
\newblock QLoRA: Efficient finetuning of quantized LLMs.
\newblock En \emph{Advances in Neural Information Processing Systems (NeurIPS)}, 2023.
\newblock \url{https://arxiv.org/abs/2305.14314}

\bibitem{eckart-young}
C.~Eckart y G.~Young.
\newblock The approximation of one matrix by another of lower rank.
\newblock \emph{Psychometrika}, 1(3):211--218, 1936.

\bibitem{gemini}
Gemini Team, Google.
\newblock Gemini: A family of highly capable multimodal models.
\newblock Informe técnico, 2023.
\newblock \url{https://arxiv.org/abs/2312.11805}

\bibitem{gbc}
I.~Goodfellow, Y.~Bengio y A.~Courville.
\newblock \emph{Deep Learning}.
\newblock MIT Press, 2016. §7.1.1, «$L^2$ Parameter Regularization».
\newblock \url{https://www.deeplearningbook.org/}

\bibitem{he-unified}
J.~He, C.~Zhou, X.~Ma, T.~Berg-Kirkpatrick y G.~Neubig.
\newblock Towards a unified view of parameter-efficient transfer learning.
\newblock En \emph{International Conference on Learning Representations (ICLR)}, 2022.
\newblock \url{https://arxiv.org/abs/2110.04366}

\bibitem{houlsby}
N.~Houlsby, A.~Giurgiu, S.~Jastrzebski, B.~Morrone, Q.~de~Laroussilhe, A.~Gesmundo, M.~Attariyan y S.~Gelly.
\newblock Parameter-efficient transfer learning for NLP.
\newblock En \emph{International Conference on Machine Learning (ICML)}, 2019.
\newblock \url{https://arxiv.org/abs/1902.00751}

\bibitem{howard-ruder}
J.~Howard y S.~Ruder.
\newblock Universal language model fine-tuning for text classification.
\newblock En \emph{Proceedings of ACL}, 2018.
\newblock \url{https://arxiv.org/abs/1801.06146}

\bibitem{hu-lora}
E.~J. Hu, Y.~Shen, P.~Wallis, Z.~Allen-Zhu, Y.~Li, S.~Wang, L.~Wang y W.~Chen.
\newblock LoRA: Low-rank adaptation of large language models.
\newblock En \emph{International Conference on Learning Representations (ICLR)}, 2022.
\newblock \url{https://arxiv.org/abs/2106.09685}

\bibitem{jm3}
D.~Jurafsky y J.~H. Martin.
\newblock \emph{Speech and Language Processing}, 3.ª ed., borrador del 19 de agosto de 2026.
\newblock Cap.~8, «Post-training»; cap.~11, «Information Retrieval and RAG».
\newblock \url{https://web.stanford.edu/~jurafsky/slp3/}

\bibitem{kirkpatrick}
J.~Kirkpatrick \emph{et al.}
\newblock Overcoming catastrophic forgetting in neural networks.
\newblock \emph{Proceedings of the National Academy of Sciences}, 114(13):3521--3526, 2017.
\newblock \url{https://arxiv.org/abs/1612.00796}

\bibitem{lester}
B.~Lester, R.~Al-Rfou y N.~Constant.
\newblock The power of scale for parameter-efficient prompt tuning.
\newblock En \emph{Proceedings of EMNLP}, 2021.
\newblock \url{https://arxiv.org/abs/2104.08691}

\bibitem{li-intrinsic}
C.~Li, H.~Farkhoor, R.~Liu y J.~Yosinski.
\newblock Measuring the intrinsic dimension of objective landscapes.
\newblock En \emph{International Conference on Learning Representations (ICLR)}, 2018.
\newblock \url{https://arxiv.org/abs/1804.08838}

\bibitem{li-liang}
X.~L. Li y P.~Liang.
\newblock Prefix-tuning: Optimizing continuous prompts for generation.
\newblock En \emph{Proceedings of ACL-IJCNLP}, 2021.
\newblock \url{https://arxiv.org/abs/2101.00190}

\bibitem{mccloskey}
M.~McCloskey y N.~J. Cohen.
\newblock Catastrophic interference in connectionist networks: The sequential learning problem.
\newblock \emph{Psychology of Learning and Motivation}, 24:109--165, 1989.

\bibitem{ouyang}
L.~Ouyang \emph{et al.}
\newblock Training language models to follow instructions with human feedback.
\newblock En \emph{Advances in Neural Information Processing Systems (NeurIPS)}, 2022.
\newblock \url{https://arxiv.org/abs/2203.02155}

\bibitem{sanh-t0}
V.~Sanh \emph{et al.}
\newblock Multitask prompted training enables zero-shot task generalization.
\newblock En \emph{International Conference on Learning Representations (ICLR)}, 2022.
\newblock \url{https://arxiv.org/abs/2110.08207}

\bibitem{wei-flan}
J.~Wei, M.~Bosma, V.~Y. Zhao, K.~Guu, A.~W. Yu, B.~Lester, N.~Du, A.~M. Dai y Q.~V. Le.
\newblock Finetuned language models are zero-shot learners.
\newblock En \emph{International Conference on Learning Representations (ICLR)}, 2022.
\newblock \url{https://arxiv.org/abs/2109.01652}

\bibitem{d2l}
A.~Zhang, Z.~C. Lipton, M.~Li y A.~J. Smola.
\newblock \emph{Dive into Deep Learning}, versión 1.0.3.
\newblock Cambridge University Press, 2023. §14.2, «Fine-Tuning».
\newblock \url{https://d2l.ai/}

\bibitem{zhou-lima}
C.~Zhou \emph{et al.}
\newblock LIMA: Less is more for alignment.
\newblock En \emph{Advances in Neural Information Processing Systems (NeurIPS)}, 2023.
\newblock \url{https://arxiv.org/abs/2305.11206}

\end{thebibliography}


\end{document}
